% GENERATED FILE. Edit canonical lessons, then run npm run generate:books.
% canonical-source-hash: fnv1a64:c904f801fe60d9de
% canonical-lessons: AR-C330-sajana, AR-C330-sajin, AR-C330-zinzana, AR-C330-qayd, AR-C330-afraja

\chapter{Prison}
\label{ch:ar-prison}
\hlchaptermodality{330}

\begin{chapteropening}
\textbf{By the end of this chapter:} \emph{I can say to imprison, a prisoner, a cell, a shackle and to release.}

Complete the last lesson of chapter 330: I can say to imprison, a prisoner, a cell, a shackle and to release.
\end{chapteropening}

\section[\texorpdfstring{\ar{سجن} (sajana)}{سجن (sajana)}]{\ar{سجن} (sajana) — to imprison, to jail}

\label{lesson:AR-C330-sajana}



\begin{quote}
Before the new one: say the Arabic for independence, then the Arabic for a country’s people.
\end{quote}

\subsection*{You'll want to know: \ar{سجن}}
\textbf{\ar{سجن}} — \emph{sajana} — \textquotedblleft{}to imprison, to jail\textquotedblright{}.

\begin{quote}
\ar{سجن} \ar{اللص} \ar{سنتين}.
\end{quote}

\begin{quote}
\emph{sujina al-liṣṣ sanatayn.}
\end{quote}

\begin{quote}
The thief was jailed for two years.
\end{quote}

Present: \ar{يسجن}, \emph{yasjun}.

The noun \ar{سجن}, \emph{sijn}, a prison, is spelled the same.

\begin{cousinweb}
\textbf{\ar{سجن}} (\emph{sajana}) is built on the native Arabic root \emph{s-j-n} (Lane), \textquotedblleft{}to shut up, to confine\textquotedblright{}; \emph{sijn} is the place of it.
\end{cousinweb}

\begin{grammarlens}[title={What you've built}]
One more word for telling what happened.
\end{grammarlens}

\subsection*{Your turn}
\begin{itemize}
\raggedright
  \item \emph{Say it:} \emph{sajana}
  \item \emph{Say it:} \emph{sajana}, once more
  \item \emph{Say it:} \emph{sujina al-liṣṣ sanatayn.} — The thief was jailed for two years.
  \item \emph{Recall:} say the Arabic for independence, then the Arabic for a country’s people, then say \emph{sajana} again
\end{itemize}

\begin{tcolorbox}[breakable,colback=teal!4,colframe=teal!35!black,title={Before you move on}]
What does \ar{سجن} mean? (\textquotedblleft{}To imprison, to jail\textquotedblright{}.) Say it once more.
\end{tcolorbox}
\section[\texorpdfstring{\ar{سجين} (sajīn)}{سجين (sajīn)}]{\ar{سجين} (sajīn) — a prisoner}

\label{lesson:AR-C330-sajin}



\begin{quote}
Before the new one: say the Arabic for a country’s people, then the Arabic for to imprison.
\end{quote}

\subsection*{You'll want to know: \ar{سجين}}
\textbf{\ar{سجين}} — \emph{sajīn} — \textquotedblleft{}a prisoner\textquotedblright{}.

\begin{quote}
\ar{هرب} \ar{سجين} \ar{من} \ar{السجن} \ar{ليلا}.
\end{quote}

\begin{quote}
\emph{haraba sajīn min as-sijn laylan.}
\end{quote}

\begin{quote}
A prisoner escaped from the prison at night.
\end{quote}

Plural: \ar{سجناء}, \emph{sujanāʾ}.

\begin{cousinweb}
\textbf{\ar{سجين}} (\emph{sajīn}) is built on the native Arabic root \emph{s-j-n} (Lane), \textquotedblleft{}to confine\textquotedblright{}. The \emph{faʿīl} pattern gives the one confined.
\end{cousinweb}

\begin{grammarlens}[title={What you've built}]
One more word for telling what happened.
\end{grammarlens}

\subsection*{Your turn}
\begin{itemize}
\raggedright
  \item \emph{Say it:} \emph{sajīn}
  \item \emph{Say it:} \emph{sajīn}, once more
  \item \emph{Say it:} \emph{haraba sajīn min as-sijn laylan.} — A prisoner escaped from the prison at night.
  \item \emph{Recall:} say the Arabic for a country’s people, then the Arabic for to imprison, then say \emph{sajīn} again
\end{itemize}

\begin{tcolorbox}[breakable,colback=teal!4,colframe=teal!35!black,title={Before you move on}]
What does \ar{سجين} mean? (\textquotedblleft{}A prisoner\textquotedblright{}.) Say it once more.
\end{tcolorbox}
\section[\texorpdfstring{\ar{زنزانة} (zinzāna)}{زنزانة (zinzāna)}]{\ar{زنزانة} (zinzāna) — a prison cell}

\label{lesson:AR-C330-zinzana}



\begin{quote}
Before the new one: say the Arabic for to imprison, then the Arabic for a prisoner.
\end{quote}

\subsection*{You'll want to know: \ar{زنزانة}}
\textbf{\ar{زنزانة}} — \emph{zinzāna} — \textquotedblleft{}a prison cell\textquotedblright{}.

\begin{quote}
\ar{كانت} \ar{الزنزانة} \ar{صغيرة} \ar{ومظلمة}.
\end{quote}

\begin{quote}
\emph{kānat az-zinzāna ṣaghīra wa-muẓlima.}
\end{quote}

\begin{quote}
The cell was small and dark.
\end{quote}

Plural: \ar{زنازين}, \emph{zanāzīn}.

\begin{cousinweb}
\textbf{\ar{زنزانة}} (\emph{zinzāna}) is not in the old dictionaries Lane drew on; it belongs to later Arabic, and the origin is not settled.
\end{cousinweb}

\begin{grammarlens}[title={What you've built}]
One more word for telling what happened.
\end{grammarlens}

\subsection*{Your turn}
\begin{itemize}
\raggedright
  \item \emph{Say it:} \emph{zinzāna}
  \item \emph{Say it:} \emph{zinzāna}, once more
  \item \emph{Say it:} \emph{kānat az-zinzāna ṣaghīra wa-muẓlima.} — The cell was small and dark.
  \item \emph{Recall:} say the Arabic for to imprison, then the Arabic for a prisoner, then say \emph{zinzāna} again
\end{itemize}

\begin{tcolorbox}[breakable,colback=teal!4,colframe=teal!35!black,title={Before you move on}]
What does \ar{زنزانة} mean? (\textquotedblleft{}A prison cell\textquotedblright{}.) Say it once more.
\end{tcolorbox}
\section[\texorpdfstring{\ar{قيد} (qayd)}{قيد (qayd)}]{\ar{قيد} (qayd) — a shackle; a restriction}

\label{lesson:AR-C330-qayd}



\begin{quote}
Before the new one: say the Arabic for a prisoner, then the Arabic for a prison cell.
\end{quote}

\subsection*{You'll want to know: \ar{قيد}}
\textbf{\ar{قيد}} — \emph{qayd} — \textquotedblleft{}a shackle; a restriction\textquotedblright{}.

\begin{quote}
\ar{وضعت} \ar{الشرطة} \ar{القيد} \ar{في} \ar{يدي} \ar{اللص}.
\end{quote}

\begin{quote}
\emph{waḍaʿat ash-shurṭa al-qayd fī yaday al-liṣṣ.}
\end{quote}

\begin{quote}
The police put the handcuffs on the thief’s hands.
\end{quote}

Plural: \ar{قيود}, \emph{quyūd}.

\ar{على} \ar{قيد} \ar{الحياة}: alive, \textquotedblleft{}in the bond of life\textquotedblright{}.

\begin{cousinweb}
\textbf{\ar{قيد}} (\emph{qayd}) is built on the native Arabic root \emph{q-y-d} (Lane), where \emph{qayd} is the hobble put on a horse or camel, and \emph{qayyada}, \textquotedblleft{}to tie; to note down\textquotedblright{}.
\end{cousinweb}

\begin{grammarlens}[title={What you've built}]
One more word for telling what happened.
\end{grammarlens}

\subsection*{Your turn}
\begin{itemize}
\raggedright
  \item \emph{Say it:} \emph{qayd}
  \item \emph{Say it:} \emph{qayd}, once more
  \item \emph{Say it:} \emph{waḍaʿat ash-shurṭa al-qayd fī yaday al-liṣṣ.} — The police put the handcuffs on the thief’s hands.
  \item \emph{Recall:} say the Arabic for a prisoner, then the Arabic for a prison cell, then say \emph{qayd} again
\end{itemize}

\begin{tcolorbox}[breakable,colback=teal!4,colframe=teal!35!black,title={Before you move on}]
What does \ar{قيد} mean? (\textquotedblleft{}A shackle; a restriction\textquotedblright{}.) Say it once more.
\end{tcolorbox}
\section[\texorpdfstring{\ar{أفرج} (afraja)}{أفرج (afraja)}]{\ar{أفرج} (afraja) — to release, to set free}

\label{lesson:AR-C330-afraja}



\begin{quote}
Before the new one: say the Arabic for a prison cell, then the Arabic for a shackle.
\end{quote}

\subsection*{You'll want to know: \ar{أفرج}}
\textbf{\ar{أفرج}} — \emph{afraja} — \textquotedblleft{}to release, to set free\textquotedblright{}.

\begin{quote}
\ar{أفرجت} \ar{الشرطة} \ar{عن} \ar{الرجل} \ar{في} \ar{المساء}.
\end{quote}

\begin{quote}
\emph{afrajat ash-shurṭa ʿan ar-rajul fī al-masāʾ.}
\end{quote}

\begin{quote}
The police released the man in the evening.
\end{quote}

Present: \ar{يفرج}, \emph{yufrij}.

It takes \ar{عن}: \ar{أفرج} \ar{عن} \ar{السجين}.

\begin{cousinweb}
\textbf{\ar{أفرج}} (\emph{afraja}) is built on the native Arabic root \emph{f-r-j} (Lane), \textquotedblleft{}to open a gap; to relieve\textquotedblright{}; \emph{faraj} is relief from trouble. It is Form IV, \textquotedblleft{}to open the way out\textquotedblright{}.
\end{cousinweb}

\begin{grammarlens}[title={What you've built}]
One more word for telling what happened.
\end{grammarlens}

\subsection*{Your turn}
\begin{itemize}
\raggedright
  \item \emph{Say it:} \emph{afraja}
  \item \emph{Say it:} \emph{afraja}, once more
  \item \emph{Say it:} \emph{afrajat ash-shurṭa ʿan ar-rajul fī al-masāʾ.} — The police released the man in the evening.
  \item \emph{Recall:} say the Arabic for a prison cell, then the Arabic for a shackle, then say \emph{afraja} again
\end{itemize}

\begin{tcolorbox}[breakable,colback=teal!4,colframe=teal!35!black,title={Before you move on}]
What does \ar{أفرج} mean? (\textquotedblleft{}To release, to set free\textquotedblright{}.) Say it once more.
\end{tcolorbox}
